\addtocontents{toc}{\protect\thispagestyle{empty}} % força o "empty" na primeira pagina do sumário
\pagestyle{empty} % remove a numeracao das demais paginas de sumario (se houver mais de uma) 

%apresenta ate o nivel 5, de paragrafo
\setcounter{tocdepth}{4}

%apresenta a numeracao ate o nivel 5, de paragrafo
\setcounter{secnumdepth}{4}

% A definição desse comando que estabelece qual o ponto de início do texto (ou seja a margem do texto) dos ítens do sumário.
% Ajuste o valor do espaçamento para a que representa o nível de maior profundidade do seu trabalho
% Por exemplo, se o seu trabalho possui uma subsecao de nível x.x.x (ex 1.2.3), deixe descomentado apenas a configuração de terceiro nível, de valor 32pt. O padrão desse documento são dois níveis (valor 23pt)
% 1 - nível: \newcommand{\sumariotextmarginlenght}{13pt}
% 2 - nível: \newcommand{\sumariotextmarginlenght}{23pt}
% 3 - nível: \newcommand{\sumariotextmarginlenght}{32pt}
% 4 - nível: \newcommand{\sumariotextmarginlenght}{38pt}
% 5 - níveil: \newcommand{\sumariotextmarginlenght}{46pt}
\newcommand{\sumariotextmarginlenght}{23pt}

% titletoc - Configuração da formatação dos títulos (em seus vários níveis) no sumário. 
% A formatação desses mesmos títulos dentro do texto, precisa ser feito antes do begin{document} e ficou no arquivo: config/configuracoes.tex
\titlecontents{chapter}[\sumariotextmarginlenght]
{\normalfont\normalsize\bfseries}
{\contentslabel{\sumariotextmarginlenght}\uppercase} % removido \MakeUppercase
{}
{\titlerule*[0.4pc]{.}\contentspage} % aqui sim!

\titlecontents{section}[\sumariotextmarginlenght]
{\normalfont\normalsize}
{\contentslabel{\sumariotextmarginlenght}\uppercase}
{}
{\titlerule*[0.4pc]{.}\contentspage}

\titlecontents{subsection}[\sumariotextmarginlenght]
{\normalfont\normalsize\bfseries}
{\contentslabel{\sumariotextmarginlenght}}
{}
{\titlerule*[0.4pc]{.}\contentspage}

\titlecontents{subsubsection}[\sumariotextmarginlenght]
{\normalfont\normalsize}
{\contentslabel{\sumariotextmarginlenght}}
{}
{\titlerule*[0.4pc]{.}\contentspage}

\titlecontents{paragraph}[\sumariotextmarginlenght]
{\normalfont\normalsize}
{\contentslabel{\sumariotextmarginlenght}}
{}
{\titlerule*[0.4pc]{.}\contentspage}


%Altera e Formata o titulo do sumario
\renewcommand{\contentsname}{\normalsize\centerline{\bfseries\uppercase{Sumário}}}

\tableofcontents % Insere o sumário
\clearpage % Garante que o sumário termine aqui
\pagestyle{fancy} % Retorna ao estilo "fancy" para o restante do documents